% ============================================================
%  Übungsaufgaben Template (my_dhbw Stil, uebung_*.tex)
%  Header "Übungsblatt", grün eingefärbte <details>-Lösungen.
%  Renderer: pdflatex (zweimal)
% ============================================================
\documentclass[11pt, a4paper]{article}

\usepackage[ngerman]{babel}
\usepackage{fontspec}
\setmainfont[
  Path=/usr/share/fonts/TTF/,
  Extension=.ttf,
  BoldFont=DejaVuSans-Bold,
  ItalicFont=DejaVuSans-Oblique,
  BoldItalicFont=DejaVuSans-BoldOblique
]{DejaVuSans}
\setsansfont[
  Path=/usr/share/fonts/TTF/,
  Extension=.ttf,
  BoldFont=DejaVuSans-Bold,
  ItalicFont=DejaVuSans-Oblique,
  BoldItalicFont=DejaVuSans-BoldOblique
]{DejaVuSans}
\setmonofont[Path=/usr/share/fonts/TTF/,Extension=.ttf]{DejaVuSansMono}
\usepackage[top=2cm, bottom=2cm, left=2.4cm, right=2.4cm, footskip=1cm]{geometry}
\usepackage{amsmath, amssymb}
\usepackage{xcolor}
\usepackage{enumitem}
\usepackage{fancyhdr}
\usepackage{lastpage}
\usepackage{tabularx}
\usepackage{booktabs}
\usepackage{microtype}
\usepackage{parskip}

\usepackage{array}
\usepackage{longtable}
\usepackage{ltablex}
\keepXColumns
\usepackage{colortbl}
\usepackage[most,breakable]{tcolorbox}
\usepackage{listings}
\usepackage[normalem]{ulem}
\usepackage{pdflscape}
\usepackage{titlesec}
\usepackage[colorlinks=true, linkcolor=accent, urlcolor=accent]{hyperref}

\renewcommand{\familydefault}{\sfdefault}

% ── Theme ─────────────────────────────────────────────────────
\definecolor{accent}{RGB}{0, 60, 113}
\definecolor{primaryblue}{RGB}{0, 60, 113}
\definecolor{loesung}{RGB}{0, 100, 0}
\definecolor{codebg}{RGB}{245, 245, 245}
\definecolor{figurebg}{RGB}{232, 241, 255}
\definecolor{tablehintbg}{RGB}{240, 255, 240}
\definecolor{solutionbg}{RGB}{240, 252, 240}
\definecolor{rulegray}{RGB}{180, 180, 180}
\definecolor{tableheadbg}{RGB}{0, 60, 113}

% ── Header/Footer ─────────────────────────────────────────────
\pagestyle{fancy}
\fancyhf{}
\renewcommand{\headrulewidth}{0.4pt}
\fancyhead[L]{\footnotesize\color{accent}\nichttitle}
\fancyhead[R]{\footnotesize\color{accent}\nichtsubtitle}
\fancyfoot[R]{\footnotesize\color{accent}Blatt \thepage\,/\,\pageref{LastPage}}

% ── Typografie ────────────────────────────────────────────────
\titleformat{\section}
    {\color{accent}\large\bfseries}{}{0pt}{}
    [\vspace{-4pt}\color{accent}\rule{\linewidth}{1.5pt}]
\titleformat{\subsection}
    {\normalsize\bfseries\color{accent!85!black}}{}{0pt}{}{}
\titleformat{\subsubsection}
    {\small\bfseries\color{accent!70!black}}{}{0pt}{}{}
\titlespacing*{\section}{0pt}{10pt}{3pt}
\titlespacing*{\subsection}{0pt}{6pt}{2pt}
\titlespacing*{\subsubsection}{0pt}{4pt}{1pt}

\setlength{\parindent}{0pt}
\setlength{\parskip}{6pt}
\renewcommand{\arraystretch}{1.2}
\setlist[itemize]{noitemsep, topsep=3pt, parsep=0pt, leftmargin=1.4em}
\setlist[enumerate]{noitemsep, topsep=3pt, parsep=0pt, leftmargin=1.6em}

% ── Boxen: solutionbox = Lösungsbox in grün ──────────────────
\newtcolorbox{figurebox}[1]{
  colback=figurebg, colframe=accent!50,
  title={Abbildung: #1}, fonttitle=\small\bfseries,
  breakable, before skip=6pt, after skip=6pt
}
\newtcolorbox{tablehintbox}[1]{
  colback=tablehintbg, colframe=green!50!black!50,
  title={Tabelle: #1}, fonttitle=\small\bfseries,
  breakable, before skip=6pt, after skip=6pt
}
\newtcolorbox{solutionbox}[1]{
  enhanced,
  colback=solutionbg, colframe=loesung,
  title={#1}, fonttitle=\small\bfseries,
  coltitle=white, colbacktitle=loesung,
  breakable, before skip=8pt, after skip=8pt
}

\lstset{
  basicstyle=\ttfamily\small,
  backgroundcolor=\color{codebg},
  breaklines=true,
  frame=single, framerule=0pt,
  xleftmargin=4pt, xrightmargin=4pt,
  aboveskip=6pt, belowskip=6pt,
  keepspaces=true
}

\providecommand{\nichttitle}{Übungsblatt}
\providecommand{\nichtsubtitle}{}

\renewcommand{\nichttitle}{Analysis 2}
\renewcommand{\nichtsubtitle}{Übungsblatt 2}


\begin{document}

\begin{center}
{\Large\bfseries\color{accent}\nichttitle}
\ifx\nichtsubtitle\empty\else
  \\[2pt]{\normalsize\color{accent!80!black}\nichtsubtitle}
\fi
\\[2pt]
\color{accent}\rule{0.4\linewidth}{0.8pt}\color{black}
\end{center}
\vspace{4pt}

% ── BODY ──────────────────────────────────────────────────────

\section{Übungsblatt 2 — Partielle Ableitungen}


\noindent\textcolor{rulegray}{\rule{\linewidth}{0.4pt}}


\paragraph{Aufgabe 1 — Begriff und Konsequenz \textit{(8 Punkte)}}\mbox{}\\[2pt]

a) Erkläre in eigenen Worten, was die partielle Ableitung ∂f/∂x einer Funktion f(x, y, z) bedeutet und wie sie konkret gebildet wird. \textit{(3 P)}


b) Sei g(x, y, z) = 7y³ + 2z² − 5. Bestimme ∂g/∂x und begründe das Ergebnis aus der Definition. Welche allgemeine Konsequenz ergibt sich für Summanden, in denen die Differenzierungsvariable nicht vorkommt? \textit{(5 P)}


\textbf{Lösung:}


a) Die partielle Ableitung ∂f/∂x ist die Ableitung der Funktion f nach der Variable x, wobei alle anderen Variablen (hier y und z) als Konstanten behandelt werden. Praktisch: man wendet die normalen Ableitungsregeln (Summen-, Produkt-, Quotienten-, Kettenregel) nur auf den x-Anteil an, alles ohne x bleibt unverändert oder fällt — falls es ein reiner konstanter Summand ohne x ist — komplett weg.


b) g(x, y, z) = 7y³ + 2z² − 5. Keiner der drei Summanden enthält die Variable x. Da y, z bei der Differentiation nach x als Konstanten gelten, sind 7y³, 2z² und −5 allesamt konstante Summanden bezüglich x. Konstante Summanden ohne Differenzierungsvariable entfallen vollständig:


∂g/∂x = 0.


\textbf{Konsequenz:} Hängt eine Funktion in einem Summanden gar nicht von der Differenzierungsvariable ab, liefert dieser Summand bei der partiellen Ableitung den Beitrag 0. Das ist auch der Grund, warum bei f(x,y) = h(y) + k(x) die partielle Ableitung nach x nur k'(x) liefert — h(y) verschwindet.



\noindent\textcolor{rulegray}{\rule{\linewidth}{0.4pt}}


\paragraph{Aufgabe 2 — Tabelle vervollständigen \textit{(15 Punkte)}}\mbox{}\\[2pt]

Vervollständige die folgende Tabelle. Gib jeweils ∂f/∂x, ∂f/∂y und (falls anwendbar) ∂f/∂z an. Zeige bei (ii) und (iii) den Rechenweg.



{\small
\begin{tabularx}{\linewidth}{XX}
\hline
\rowcolor{tableheadbg}
{\color{white}\textbf{Nr.}} & {\color{white}\textbf{f(x, y, z)}} \\[2pt]
\hline
\endhead
\hline
\endfoot
(i) & 4x⁵ − 3xy + 2y² \\
(ii) & (x² + y) / z \\
(iii) & √(x + y²) · z \\
\end{tabularx}
}


\textit{(je 5 P)}


\textbf{Lösung:}


\textbf{(i) f = 4x⁵ − 3xy + 2y²}


\begin{itemize}
  \item ∂f/∂x: 4x⁵ → 20x⁴; −3xy → −3y (y konstant); 2y² → 0. Also \textbf{∂f/∂x = 20x⁴ − 3y}.
  \item ∂f/∂y: 4x⁵ → 0; −3xy → −3x; 2y² → 4y. Also \textbf{∂f/∂y = −3x + 4y}.
  \item ∂f/∂z = 0 (z kommt nicht vor).
\end{itemize}

\textbf{(ii) f = (x² + y) / z}


Schreibe als f = (x² + y) · z⁻¹.


\begin{itemize}
  \item ∂f/∂x: y und z konstant. Ableitung von x² ist 2x, also \textbf{∂f/∂x = 2x/z}.
  \item ∂f/∂y: x und z konstant. Ableitung von y ist 1, also \textbf{∂f/∂y = 1/z}.
  \item ∂f/∂z: (x² + y) konstant, ∂/∂z (z⁻¹) = −z⁻² = −1/z². Also \textbf{∂f/∂z = −(x² + y)/z²}.
\end{itemize}

\textbf{(iii) f = √(x + y²) · z = (x + y²)\textasciicircum{}\{0,5\} · z}


Innere Ableitung mit Kettenregel.


\begin{itemize}
  \item ∂f/∂x: z konstant; ∂/∂x (x + y²)\textasciicircum{}\{0,5\} = 0,5·(x + y²)\textasciicircum{}\{−0,5\} · 1 = 1/(2√(x + y²)). Also \textbf{∂f/∂x = z / (2√(x + y²))}.
  \item ∂f/∂y: z konstant; ∂/∂y (x + y²)\textasciicircum{}\{0,5\} = 0,5·(x + y²)\textasciicircum{}\{−0,5\} · 2y = y/√(x + y²). Also \textbf{∂f/∂y = yz / √(x + y²)}.
  \item ∂f/∂z: (x + y²)\textasciicircum{}\{0,5\} konstant; ∂/∂z (z) = 1. Also \textbf{∂f/∂z = √(x + y²)}.
\end{itemize}


\noindent\textcolor{rulegray}{\rule{\linewidth}{0.4pt}}


\paragraph{Aufgabe 3 — Anwendung mit Logarithmus \textit{(12 Punkte)}}\mbox{}\\[2pt]

Gegeben sei


f(x, y, z) = 3x · ln(x²yz)


a) Bestimme ∂f/∂x mit Produkt- und Kettenregel. \textit{(5 P)}


b) Bestimme ∂f/∂y und ∂f/∂z. \textit{(4 P)}


c) Werte alle drei partiellen Ableitungen am Punkt (x, y, z) = (1, e, 1) aus. \textit{(3 P)}


\textbf{Lösung:}


\textbf{a) ∂f/∂x:} Produktregel auf 3x · ln(x²yz).


\begin{itemize}
  \item (3x)' = 3.
  \item ln(x²yz)' nach x: innere Ableitung ∂/∂x (x²yz) = 2xyz, also ∂/∂x ln(x²yz) = 2xyz / (x²yz) = 2/x.
\end{itemize}

Damit:

∂f/∂x = 3 · ln(x²yz) + 3x · (2/x) = \textbf{3 ln(x²yz) + 6}.


\textbf{b)}


∂f/∂y: 3x ist bzgl. y konstant.

∂/∂y ln(x²yz) = (x²z)/(x²yz) = 1/y.

Also ∂f/∂y = 3x · (1/y) = \textbf{3x/y}.


∂f/∂z: analog ∂/∂z ln(x²yz) = (x²y)/(x²yz) = 1/z.

Also ∂f/∂z = 3x · (1/z) = \textbf{3x/z}.


\textbf{c) Auswertung in (1, e, 1):}


x²yz = 1 · e · 1 = e, also ln(x²yz) = ln(e) = 1.


\begin{itemize}
  \item ∂f/∂x = 3·1 + 6 = \textbf{9}.
  \item ∂f/∂y = 3·1/e = \textbf{3/e}.
  \item ∂f/∂z = 3·1/1 = \textbf{3}.
\end{itemize}


\noindent\textcolor{rulegray}{\rule{\linewidth}{0.4pt}}


\paragraph{Aufgabe 4 — Transfer: Wärmeverteilung an einer Platte \textit{(20 Punkte)}}\mbox{}\\[2pt]

Auf einer rechteckigen Metallplatte ist die Temperatur (in °C) modelliert durch


T(x, y) = 60 + 4x² − xy + 2y²,


wobei x, y die Koordinaten in Metern sind. Ein Sensor sitzt am Punkt P = (3, 2).


a) Bestimme ∂T/∂x und ∂T/∂y. \textit{(4 P)}


b) Werte beide partielle Ableitungen am Punkt P aus und interpretiere die Vorzeichen und Beträge physikalisch (welche Richtung bringt schnellere Temperaturänderung?). \textit{(6 P)}


c) Ein zweiter Sensor sitzt bei Q = (1, 5). In welche Achsen­richtung (x oder y) ist die Temperatur dort lokal empfindlicher gegen Verschiebung? Begründe rechnerisch. \textit{(5 P)}


d) An welchen Punkten der Platte gilt simultan ∂T/∂x = 0 und ∂T/∂y = 0? Bestimme alle Lösungen. \textit{(5 P)}


\textbf{Lösung:}


\textbf{a)} Konstante 60 fällt weg.


∂T/∂x = 8x − y.

∂T/∂y = −x + 4y.


\textbf{b)} In P = (3, 2):


∂T/∂x|\_P = 8·3 − 2 = 22 [°C/m].

∂T/∂y|\_P = −3 + 4·2 = 5 [°C/m].


Interpretation: Bei einer kleinen Verschiebung in x-Richtung wächst die Temperatur um ca. 22 °C pro Meter, in y-Richtung um ca. 5 °C pro Meter. Beide Vorzeichen sind positiv → in beide Achsenrichtungen wird es wärmer. Die x-Richtung ist deutlich „heißempfindlicher": ∂T/∂x ist mehr als viermal so groß wie ∂T/∂y.


\textbf{c)} In Q = (1, 5):


∂T/∂x|\_Q = 8·1 − 5 = 3.

∂T/∂y|\_Q = −1 + 4·5 = 19.


|∂T/∂y| = 19 > |∂T/∂x| = 3 → in Q ist die Temperatur \textbf{in y-Richtung} lokal empfindlicher: eine Verschiebung um Δy bewirkt eine ca. 6,3-fach stärkere Temperaturänderung als die gleiche Verschiebung in x.


\textbf{d)} System:


8x − y = 0 → y = 8x.

−x + 4y = 0 → x = 4y.


Einsetzen: x = 4·(8x) = 32x → 31x = 0 → x = 0, daraus y = 0.


Einzige Lösung: \textbf{(x, y) = (0, 0)}. An diesem Punkt verschwinden beide partiellen Ableitungen gleichzeitig (kandidatischer kritischer Punkt der Funktion).



\noindent\textcolor{rulegray}{\rule{\linewidth}{0.4pt}}


\paragraph{Aufgabe 5 — Fehler finden \textit{(10 Punkte)}}\mbox{}\\[2pt]

Eine Kommilitonin reicht die folgende Rechnung ein. Gegeben sei


f(x, y, z) = (x + y²) / z + 2xy.


Ihre Behauptung:


\begin{quote}
∂f/∂y = 2y/z + 2x − (x + y²)/z².
\end{quote}

Finde \textbf{alle} Fehler, korrigiere sie und gib die richtige partielle Ableitung ∂f/∂y vollständig an. Erkläre kurz, woher der Fehler vermutlich kommt.


\textbf{Lösung:}


Wir leiten Term für Term nach y ab. Bei ∂/∂y werden x und z als Konstanten behandelt.


\begin{itemize}
  \item Term 1: (x + y²)/z. Nur y² hängt von y ab. ∂/∂y ((x + y²)/z) = 2y/z. ✓
  \item Term 2: 2xy. ∂/∂y (2xy) = 2x. ✓
\end{itemize}

Also: \textbf{∂f/∂y = 2y/z + 2x}.


\textbf{Fehler:} Die Kommilitonin hat zusätzlich den Term −(x + y²)/z² angehängt. Dieser Term ist die Ableitung von (x + y²)/z \textbf{nach z}, nicht nach y. Vermutlich hat sie die Spalte ∂f/∂z aus der Beispieltabelle versehentlich übernommen (klassischer Spalten-Verwechsler) oder beim Umschreiben (x + y²)·z⁻¹ einmal nach z statt nach y abgeleitet.


Der Term −(x + y²)/z² gehört korrekt nur in ∂f/∂z = −(x + y²)/z² + 0 = −(x + y²)/z².



\noindent\textcolor{rulegray}{\rule{\linewidth}{0.4pt}}


\paragraph{Aufgabe 6 — Vergleich zweier Schreibweisen \textit{(15 Punkte)}}\mbox{}\\[2pt]

Gegeben sei


f(x, y, z) = √(x·y) / z.


Zwei Studenten lösen die Aufgabe ∂f/∂x unterschiedlich:


\begin{itemize}
  \item \textbf{Variante A} schreibt f = (xy)\textasciicircum{}\{0,5\} · z⁻¹ und differenziert direkt nach x.
  \item \textbf{Variante B} schreibt f = √x · √y / z und nutzt, dass √y/z bzgl. x konstant ist.
\end{itemize}

a) Führe beide Varianten vollständig durch und zeige, dass das Ergebnis identisch ist. \textit{(8 P)}


b) Welche Variante ist hier rechnerisch effizienter? Begründe und nenne ein Szenario (eigenes Beispiel mit anderer Funktion), in dem die jeweils andere Variante besser ist. \textit{(7 P)}


\textbf{Lösung:}


\textbf{a)}


\textbf{Variante A:} f = (xy)\textasciicircum{}\{0,5\} · z⁻¹. Nach x ableiten, z⁻¹ ist konstant bzgl. x:


∂f/∂x = z⁻¹ · 0,5·(xy)\textasciicircum{}\{−0,5\} · y = (0,5 y) / (z·√(xy)) = y / (2z√(xy)).


Vereinfachen: √(xy) = √x·√y, also y/(2z√x·√y) = √y / (2z√x).


\textbf{Ergebnis A:} ∂f/∂x = √y / (2z√x).


\textbf{Variante B:} f = √x · (√y / z). Faktor √y/z ist bzgl. x konstant:


∂f/∂x = (√y/z) · ∂/∂x (√x) = (√y/z) · 1/(2√x) = √y / (2z√x).


\textbf{Ergebnis B:} ∂f/∂x = √y / (2z√x). ✓ identisch.


\textbf{b) Vergleich:}


Variante B ist hier effizienter: durch das Faktorisieren √x · (√y/z) wird der konstante Faktor sofort ausgeklammert, und es bleibt eine einfache eindimensionale Ableitung d/dx √x = 1/(2√x). Variante A muss die Kettenregel auf (xy)\textasciicircum{}\{0,5\} anwenden und am Ende noch den Wurzelterm zerlegen, um auf die gleiche Form zu kommen — mehr Schritte, mehr Fehlerquellen.


\textbf{Szenario, in dem A besser ist:} Wenn das Argument der Wurzel nicht in x- und y-Faktoren zerfällt, z. B.


g(x, y, z) = √(x + y²) / z.


Hier kann man \textbf{nicht} in einen reinen x-Faktor und einen reinen y-Faktor trennen — die Summe x + y² lässt sich nicht zerlegen. Dann muss man die Form (x + y²)\textasciicircum{}\{0,5\} · z⁻¹ stehen lassen und mit Kettenregel direkt ableiten (Variante A): ∂g/∂x = 1/(2z√(x+y²)). Variante B ist hier nicht anwendbar.


\textbf{Trade-off:} B nur bei multiplikativer Trennbarkeit der Variablen; A funktioniert immer, ist aber bei trennbaren Funktionen unnötig aufwendig.





% ──────────────────────────────────────────────────────────────

\end{document}
